% =====================================================================
%  04-validador-evidencia.tex: Validador de paquetes de evidencia
% =====================================================================

\subsection{Validador de paquetes de evidencia}\label{sec:impl-validador}

Esta subsección implementa el diseño de la \cref{sec:validacion-navegador} (ADR~0085). Tiene tres partes: un motor de verificación puro, su ejecución dentro del navegador y la generación del manifiesto sellado durante la aceptación de una firma individual. El motor es el mismo en los dos extremos: el navegador lo usa para leer paquetes y el servidor, para escribir el manifiesto y el recibo.

\paragraph{Motor puro.} El crate \texttt{evidence-verify} no depende de ningún otro crate del workspace ni usa archivos, red, reloj o números aleatorios; por eso compila igual para el servidor y para WebAssembly. Recorre la evidencia en un orden fijo. Primero comprueba que el ZIP tenga exactamente la forma que escribe Qadra; después reconoce el tipo de paquete y revisa cada pieza: firmas, certificados, listas de revocación, sellos de tiempo \cite{adams2001,housley2009}, declaraciones y, en la versión 2, el manifiesto. Cada pieza se evalúa en su propio instante, como se explica en la \cref{sec:validacion-navegador}. El resultado es un informe estructurado con el modelo de la \cref{fig:cripto-validador-resultados}. Las reglas de bytes y de perfiles se reimplementaron a partir del código del servidor y se contrastan con él mediante pruebas de paridad sobre los mismos casos.

\paragraph{Ejecución en el navegador.} El crate \texttt{evidence-verify-wasm} envuelve el motor para WebAssembly y lo expone a un worker dedicado. La \cref{fig:impl-validador-secuencia} muestra el recorrido. La pestaña calcula la huella del archivo con la criptografía del navegador y consulta la ruta de confianza, de solo lectura, disponible para cualquier sesión y sin expediente, porque un archivo local no pertenece a ninguno. Con las raíces fijas, la respuesta de esa ruta, la política de sello admitida y el instante del reloj del equipo arma la política, y el cliente crea un worker para esa validación. El panel muestra cuatro etapas: las dos primeras (\emph{Leyendo el archivo} y \emph{Consultando las raíces de confianza del servidor}) ocurren en la pestaña; las dos últimas (\emph{Preparando el validador} y \emph{Revisando la evidencia}) las informa el worker, que entrega el informe y se cierra. Cancelar, cambiar de documento o de versión o cerrar el panel descarta cualquier respuesta tardía. El informe puede descargarse en JSON con la versión del motor, la política aplicada, las raíces, los instantes y el resultado de cada componente; el propio panel aclara que esa descarga describe la revisión y no es una firma ni un sello nuevo.

\begin{figure}[!tb]
  \centering
  % Fuente: web/apps/qadra/src/components/EvidenceValidation.svelte,
% web/apps/qadra/src/lib/evidence-verify-client.mjs, evidence-verify-worker.mjs,
% web/apps/qadra/src/lib/evidence/policy.mjs y trust-api.mjs;
% crates/evidence-verify-wasm y crates/web/src/evidence_trust/mod.rs.
\begin{tikzpicture}[actor/.append style={minimum width=2.5cm, text width=2.3cm}]
  \Actores{per/1.4/{Persona},pan/4.7/{Pestaña\\ Evidencia},wrk/8.0/{Web Worker},mot/11.1/{Motor\\ WebAssembly},srv/14.6/{Servidor: ruta de confianza}}{-11.0}

  \Msg{1.4}{4.7}{-1.1}{\textbf{1} elige el ZIP\\ (hasta 26~MiB)}
  \Paso{4.7}{-1.3}{2.7}{Lee el archivo y calcula su SHA-256 con la criptografía del navegador}
  \Msg{4.7}{14.6}{-2.75}{\textbf{2} consulta las raíces vigentes, con la sesión}
  \Msg[msgR]{14.6}{4.7}{-3.35}{raíces y CRL que el servidor admite hoy}
  \Paso{14.6}{-3.5}{2.5}{Si falla, la validación sigue; solo cambia el estado actual}
  \Paso{4.7}{-3.5}{2.7}{\textbf{3} Arma la política: raíces fijas, confianza servida, política de sello e instante del equipo}
  \Msg{4.7}{8.0}{-5.2}{\textbf{4} bytes y política}
  \Paso{8.0}{-5.4}{2.5}{Carga el módulo WebAssembly}
  \Msg[msgR]{8.0}{4.7}{-6.4}{etapas}
  \Msg{8.0}{11.1}{-6.95}{\textbf{5} valida}
  \Paso{11.1}{-7.1}{2.6}{Revisa cada componente sin red, reloj ni azar}
  \Msg[msgR]{11.1}{8.0}{-8.2}{informe}
  \Msg[msgR]{8.0}{4.7}{-8.75}{informe; el worker se cierra}
  \Paso{4.7}{-8.9}{2.7}{\textbf{6} Muestra el resultado global y por componente; ofrece el informe en JSON}
  \Msg[msgR]{4.7}{1.4}{-10.5}{resultado}
  \Paso[pasoError]{1.4}{-6.6}{2.5}{Cancelar, cambiar de versión o cerrar el panel descarta resultados tardíos}

  % Frontera
  \draw[IPNGuinda, dashed, line width=0.9pt] (12.85,-0.5) -- (12.85,-11.1);
  \node[font=\diag\bfseries, text=IPNGuinda, rotate=90, fill=white, inner sep=1.5pt, align=center]
    at (12.85,-7.6) {el ZIP no se envía al servidor};
\end{tikzpicture}

  \caption{Validación de un paquete en el navegador. La única consulta al servidor obtiene las raíces vigentes y no lleva el paquete; si falla, la validación sigue con las raíces fijas.}
  \label{fig:impl-validador-secuencia}
\end{figure}

\paragraph{Manifiesto en la aceptación.} Las firmas individuales nuevas generan el manifiesto dentro de la misma operación que las acepta (\cref{fig:impl-manifiesto-aceptacion}), no al exportar ni en una tarea posterior. Después de verificar la firma del PDF, la vinculación y el consentimiento, el servicio obtiene el sello del documento, construye el manifiesto con los bytes finales y pide un segundo sello sobre su huella. Rechaza un sello del manifiesto anterior al del documento, porque la TSA no garantiza el orden, y verifica la CRL del emisor de la TSA que acompañará al paquete. Las consultas a la TSA ocurren fuera de la transacción de base de datos, para no mantener la base bloqueada mientras responde la TSA; si el segundo sello o la CRL fallan, no se guarda nada. No hay reintento automático: reenviar la misma operación pide sellos nuevos, y una operación ya aceptada devuelve su manifiesto sin cambios.

\begin{figure}[!tb]
  \centering
  % Fuente: crates/application/src/document_signing/ (document_service.rs,
% document_manifest.rs, manifest_port.rs), crates/application/src/documents/signing_manifest.rs,
% migrations/0049_individual_signing_manifests.sql y ADR 0085.
\begin{tikzpicture}[actor/.append style={minimum width=2.5cm, text width=2.3cm}]
  \Actores{tit/1.45/{Titular\\ en Qadra},api/5.4/{API: aceptación},tsa/10.0/{TSA interna},pg/14.5/{PostgreSQL}}{-12.3}

  \Msg{1.45}{5.4}{-1.1}{\textbf{1} firmas del PDF\\ y del consentimiento}
  \Paso{5.4}{-1.25}{3.3}{Verifica la firma del PDF, la vinculación y el consentimiento; revisa vigencias y reautentica; exige que la raíz de credenciales emita el certificado de la TSA}
  \Msg{5.4}{10.0}{-3.05}{\textbf{2} huella del PDF}
  \Msg[msgR]{10.0}{5.4}{-3.6}{sello del documento}
  \Paso{5.4}{-3.75}{3.3}{\textbf{3} Construye \texttt{manifiesto.json} con los bytes finales de cada pieza}
  \Msg{5.4}{10.0}{-5.0}{\textbf{4} huella del manifiesto}
  \Msg[msgR]{10.0}{5.4}{-5.55}{sello del manifiesto}
  \Paso{5.4}{-5.7}{3.3}{Rechaza un sello anterior al del documento; verifica la CRL capturada del emisor de la TSA: firma, vigencia y TSA no revocada}
  \Paso[pasoError]{10.0}{-5.85}{2.9}{Si falla el segundo sello o la CRL, no se guarda nada y no se reintenta solo}
  \Paso{5.4}{-7.45}{3.3}{\textbf{5} Repite la revisión de vigencias y la reautenticación; fija el instante mínimo de aceptación}
  \Msg{5.4}{14.5}{-8.7}{\textbf{6} recibo, manifiesto, sello y CRL}
  \Paso{14.5}{-8.85}{2.7}{Una sola transacción asigna el instante de aceptación; un disparador diferido exige un manifiesto coherente por recibo; la tabla es inmutable}
  \Msg[msgR]{14.5}{5.4}{-10.55}{confirmado}
  \Msg[msgR]{5.4}{1.45}{-11.1}{recibo}
  \Paso{1.45}{-11.25}{2.6}{Reenviar la operación pide sellos nuevos; una ya aceptada devuelve su manifiesto}

  % Fuera de la transacción
  \draw[ESCOMGris, line width=0.5pt] (11.45,-2.85) -- (11.6,-2.85) -- (11.6,-5.6) -- (11.45,-5.6);
  \node[font=\diag, text=ESCOMGris, align=left, anchor=west] at (11.7,-4.25)
    {Las consultas a la TSA\\ quedan fuera de la\\ transacción};
\end{tikzpicture}

  \caption{Aceptación de una firma individual con manifiesto sellado. Los dos sellos se piden antes de abrir la transacción; recibo y manifiesto se guardan juntos o no se guardan.}
  \label{fig:impl-manifiesto-aceptacion}
\end{figure}

La base de datos guarda cada manifiesto junto al recibo de su operación y no permite modificarlo; al confirmar la aceptación comprueba que cada recibo nuevo tenga exactamente un manifiesto que coincida con él en documento, versión y huella (migración 0049, \cref{tab:anexo-migraciones}). Al arrancar, el servidor recalcula las huellas de cada manifiesto guardado y verifica su sello. La exportación agrega las cuatro entradas de la versión 2 a partir de lo guardado, de modo que exportar dos veces produce los mismos bytes. Las firmas aceptadas antes de que se incorporara el manifiesto conservan su evidencia y se exportan como versión 1. El despliegue debe cumplir un requisito: la raíz de credenciales tiene que ser la emisora del certificado de la TSA, porque de ella se toma la CRL que cubre a la TSA; el despliegue documentado lo cumple, pues usa una sola raíz.

\paragraph{Límites.} El motor admite hasta 32 entradas, 16~MiB para el documento, 1~MiB por CRL y 25~MiB de contenido en total; el panel acepta archivos de hasta 26~MiB, margen que cubre las cabeceras del ZIP. Cada manifiesto admite un solo firmante, una sola política de sello y la TSA interna; una TSA externa, varios firmantes, una firma del servidor sobre el manifiesto, reemplazar OpenSSL en el verificador del servidor y abrir el validador fuera de la vista del documento quedan fuera de esta entrega. El validador es código nuevo con relevancia de seguridad: su confianza descansa en la integridad del programa entregado al navegador, igual que el ayudante local (\cref{sec:impl-criptografia-local}). Las pruebas están en la \cref{sec:pruebas-validador}.

\Fuentes{ADR 0085; \rutarepo{crates/evidence-verify/}, \rutarepo{crates/evidence-verify-wasm/}, \rutarepo{web/apps/qadra/src/lib/evidence/}, \rutarepo{crates/web/src/evidence_trust/mod.rs} y \rutarepo{migrations/0049_individual_signing_manifests.sql}.}
